% Standalone input fragment; requires amsmath, amssymb, booktabs,
% longtable, url or hyperref, and the master document artifact macro.
% No bibliography dependency; msref/mspageref/msnameref link to the current manuscript.
\providecommand{\leanid}[1]{\nolinkurl{#1}}

\section{Lean formalization: present coverage and remaining work}
\label{sec:review-lean}

The present Lean development does not formalize the central-charge-eight
construction or its main theorem. It checks several useful implications
inside the finite-closure and top-bound arguments. A separate, much larger
development proves existence of coefficients satisfying the Evans--Gannon
equations; its categorical reconstruction and its connection to this VOA
are not formalized. These are substantive results with different scopes.
Neither should be described as a formal verification of the entire CFT
manuscript.

This assessment reads the revised manuscript, the local Lean sources,
their saved check reports, and the installed Mathlib sources. The pinned
toolchain is Lean \leanid{v4.33.1}; the Mathlib revision in
\artifact{lean/lakefile.toml} is
\leanid{0df444a360eaa60ab8c11dca51a86af692955474}.
Paths below are relative to the repository; paths beginning
\artifact{Mathlib/} are relative to
\artifact{lean/.lake/packages/mathlib/}. Declaration names were checked
against that checkout. This is a source and scope audit, not a new
compiler run or an audit of every transitive dependency.

\subsection{What is already checked}

All five declarations below are in the namespace
\leanid{HaagerupIzumi.CFT}. The two source files are independent of
the Evans--Gannon umbrella import \leanid{HaagerupIzumi.lean}.

\begin{longtable}{@{}P{0.34\linewidth}P{0.58\linewidth}@{}}
\toprule
Declaration and source & Exact scope of the conclusion\\
\midrule\endfirsthead
\toprule Declaration and source & Exact scope of the conclusion\\
\midrule\endhead
\bottomrule\endfoot
\leanid{nilpotent_raising_saturation},
\artifact{CFT/FiniteClosure.lean}
& A submodule stable under two linear operators is the whole module,
assuming a spanning equation modulo their two images and vanishing of
every operator word of one fixed length. Joint nilpotence is an input;
the theorem does not derive it from the actual current weights.\\
\leanid{projected_residual_recovery}, same file
& From $v-Pv\in C$ and $Pv\in C$, deduce $v\in C$.
No idempotence of $P$ is needed. The actual projector coefficients and
the two membership assertions are not checked by this theorem.\\
\leanid{jacobi_bottom_witness}, same file
& From $f(v)=g(u)$, $u=0$, $v=cq$, and $c\ne0$, deduce $f(q)=0$.
The Jacobi identity, the relevant quotient, and the particular bottom
coefficient are explicit inputs, not reconstructed VOA facts.\\
\leanid{core_top_deficit_identity},
\artifact{CFT/TopBound.lean}
& The printed fourteen-variable polynomial identity, with the exact
dimension and integer-scaled trace coefficient definitions.\\
\leanid{core_top_bound}, same file
& Nonnegative multiplicities and the two stated current-character
inequalities imply the cleared-denominator $17/10$ bound. It does not
derive the fourteen tops, their multiplicities, or their stress traces.\\
\end{longtable}

The full file paths have prefix \artifact{lean/HaagerupIzumi/}.
The reports \artifact{c8_revision_20260927/lean_finite_closure_final.log}
and \artifact{c8_revision_20260927/lean_top_bound.log} record successful
checks. Their axiom reports contain only \leanid{propext},
\leanid{Quot.sound}, and, where used, \leanid{Classical.choice}.
The projected-residual lemma has an empty axiom report. These are the
usual logical foundations of the stated Lean proofs, rather than new
CFT-specific assumptions. The rational reconstruction of the top table
in \artifact{c8_revision_20260927/check_top_bound.py} is separate from
the Lean proof of the inequality.

\subsection{Assessment of every manuscript section}

The difficulty descriptions distinguish a \emph{local conditional
formalization}, in which the surrounding mathematical structures and
theorems are supplied explicitly, from an \emph{end-to-end
formalization}, which constructs those structures and proves the
inputs. ``Focused'' means a bounded algebraic or analytic lemma with
nearby library support. ``Substantial'' means several interacting new
definitions and proofs. ``Foundational'' means that major mathematical
theories or deep external theorems still have to be formalized. These
are assessments of dependency scope, not estimates of working hours.

\begin{longtable}{@{}P{0.15\linewidth}P{0.25\linewidth}P{0.52\linewidth}@{}}
\toprule
Current anchor & Existing CFT coverage & Missing work and conversion difficulty\\
\midrule\endfirsthead
\toprule
Current anchor & Existing CFT coverage & Missing work and conversion difficulty\\
\midrule\endhead
\bottomrule\endfoot
Section~\msref{sec:statement}\newline p.~\mspageref{sec:statement}
& No formal main theorem or definition of the constructed $U$.
& A theorem assembling explicit interfaces would be focused once its
inputs exist. A proof of the stated existence theorem inherits all
the foundational dependencies below; packaging them as assumptions
does not discharge them.\\
Appendix~\msref{sec:construction}\newline p.~\mspageref{sec:construction}
& No lattice VOA, cocycle-to-field construction, or embedding is
formalized here.
& Finite lattice arithmetic, the four-component coefficient algebra,
and low-weight identities are substantial but concrete targets.
Construction of the Fock space, lattice fields, PCT, and their VOA
identities is foundational.\\
Appendix~\msref{sec:finite}\newline p.~\mspageref{sec:finite}
& No filtered Wick theorem or all-weight recursion is formalized.
& A graded mode-algebra interface could support the induction,
filtration, skew-symmetry, and Wick reductions. This is a substantial
local project; deriving that interface for the actual lattice VOA
adds foundational work.\\
Appendix~\msref{sec:gate4}\newline p.~\mspageref{sec:gate4}
& Raising saturation and projected-residual recovery are checked as
abstract implications.
& Prove bounded-current-weight joint nilpotence, the singlet
projection identities, and the actual rank and residual certificates.
These are focused to substantial once the finite spaces and current
actions are represented faithfully.\\
Appendix~\msref{sec:gate5}\newline p.~\mspageref{sec:gate5}
& Only the abstract bottom-witness scalar cancellation is checked.
& Derive the quotient Jacobi relation from mode identities, justify
the charged spanning and neutral-family reductions, and verify all
nonzero coefficients. The local proof is substantial; most scalar
identities are focused.\\
Appendix~\msref{sec:properness}\newline p.~\mspageref{sec:properness}
& No theorem applies the gates to the actual subVOA.
& Given a formal recursive criterion and genuine gate inputs, the
induction, dimension comparison, and properness deduction are focused.
Without the bridge from finite data to actual modes, they remain
conditional.\\
Appendix~\msref{sec:regularity}\newline p.~\mspageref{sec:regularity}
& No embedded Potts/parafermion VOA or regularity theorem is formalized.
& Finite positive-operator inequalities and the displayed stress
identity are accessible local targets. Their application requires
conformal vectors, $C_2$ spaces, regular modules, cosets, and the
relevant rationality theorems: foundational work.\\
Appendix~\msref{sec:branching}\newline p.~\mspageref{sec:branching}
& The final fourteen-variable top-bound identity and inequality are
checked.
& Derive the core-top table and current inequalities; formalize
branching, the rank-four exclusion, rank-six based-ring argument,
formal codegrees, and center data. Conditional finite arithmetic is
focused to substantial. VOA tensor theory and fusion-category
classification remain foundational. The optional classification
route must retain its extra hypothesis.\\
Appendix~\msref{sec:characters}\newline p.~\mspageref{sec:characters}
& No theorem identifies the full character vector.
& Finite twist arithmetic and the equation $6=6+13\gamma$ are focused.
The proof needs trace characters, Zhu modularity, the specified
modular representation, finite image, the principal-part uniqueness
argument, and the Evans--Gannon fundamental matrix. Their complete
formalization is foundational, despite existing scalar modular-form
infrastructure.\\
Appendix~\msref{sec:analytic-unitarity}\newline p.~\mspageref{sec:analytic-unitarity}
& None of the analytic comparison propositions is formalized.
& Common-core reduction, dense-domain transport, analytic continuation,
and the Haag-duality squeeze are bounded local targets. Conformal
nets, DHR sectors, reconstruction, annulus integration, sewing, and
fusion positivity constitute the principal foundational burden.
Detailed distinctions appear below.\\
Appendix~\msref{sec:analytic-category}\newline p.~\mspageref{sec:analytic-category}
& No unitary ribbon identification is checked.
& The ring-level scalar arguments and cyclic-group obstruction
calculations can be isolated. A proof of the categorical conclusion
needs unitary fusion categories, the relevant Izumi classification,
Morita equivalence, and the center/ribbon comparison. General rigid
and braided categories in Mathlib supply only part of this framework.\\
Appendix~\msref{sec:analytic-full-theory}\newline p.~\mspageref{sec:analytic-full-theory}
& No Q-system extension, diagonal full theory, or $E_8$ identification
is formalized.
& Once extension correspondences and categorical duality are inputs,
transport of the chosen algebra and finite dimension checks are
substantial local targets. Proving the correspondences, the full-net
construction, and the holomorphic-VOA classification is foundational.\\
Appendix~\msref{sec:comparison}\newline p.~\mspageref{sec:comparison}
& No formal identification with another physical construction.
& Most assertions are scope statements, not independent formal
theorems. One can formalize the conditional character and extension
comparisons after the earlier interfaces exist. Neither those
interfaces nor current Lean files imply a spin-chain continuum limit
or uniqueness of all $c=8$ theories.\\
Appendix~\msref{sec:ledger}\newline p.~\mspageref{sec:ledger}
& No Lean authentication or mathematical replay of the frozen data.
& A small sound certificate checker with explicit data is a concrete,
substantial target. Its soundness must connect encoded vectors and
matrices to the actual fields. Hash agreement binds files but proves
neither their residual identities nor their VOA interpretation.\\
Appendix~\msref{sec:formalization}\newline p.~\mspageref{sec:formalization}
& The five construction declarations listed above, plus elementary
support lemmas in their two files. The current appendix also records
the two thermal lemmas assessed in the physics supplement below.
& Keep theorem names, assumptions, source hashes, toolchain, and axiom
reports synchronized. This is focused maintenance; the appendix must
not enlarge the scope of the checked declarations.\\
Appendix~\msref{sec:compilation}\newline p.~\mspageref{sec:compilation}
& Successful TeX compilation is not a Lean result.
& A reproducible build manifest can record the checked artifacts and
interfaces. Mathematical review, compiler checking, finite replay,
and PDF compilation are separate checks with separate conclusions.\\
\end{longtable}

\subsection{Finite certificates and the algebraic sections}

The new companion now supplies concrete data for this next stage:
a $120$-row short-generation graph with normalization and target
expansions, $33$ projected Gate~4 vectors and all twelve relations,
the explicit Potts and parafermion membership expressions, and the
eighty-current space's matrices and cyclic-basis recipes.  These are
in \artifact{c8_reviewability_20260927/}, with separate successful
Python coordinate, residual, and matrix checks.  They remove the
need to extract those witnesses again.  They introduce no Lean
proofs: formal parsers, checker soundness, and the interpretation
of the supplied mode recipes and matrices remain to be established.

For Appendices~\msref{sec:construction}--\msref{sec:properness}
and the ledger in Appendix~\msref{sec:ledger}, the most valuable next boundary is
between the sparse finite data and the mode calculations that they
purport to represent. A certificate checker should have a proved
interpretation map from its encoded vectors into a mathematical vector
space. It should prove that its addition, scalar multiplication, mode
actions, and selected column combinations have the stated meanings.
Only then does a zero residual certify the required identity in that
space. The further identification of this space and action with the
actual lattice fields is a distinct theorem.

The coefficient algebra uses $1,s,t,st$ with $s^2=13$ and $t^2=-3$.
One can first verify multiplication and residual identities in this
four-dimensional rational algebra. Rank arguments additionally need
a field, or an appropriately injective embedding into the intended
complex coefficient field; merely writing a quotient by the two
polynomials does not establish those properties. Sparse vectors fit
\leanid{Finsupp.linearCombination} in
\artifact{Mathlib/LinearAlgebra/Finsupp/LinearCombination.lean}.
The declarations \leanid{Matrix.rank_of_det_ne_zero},
\leanid{Matrix.rank_submatrix_le}, and
\leanid{Matrix.rank_eq_finrank_span_cols} in
\artifact{Mathlib/LinearAlgebra/Matrix/Rank.lean} provide useful rank
interfaces. A nonzero minor together with an independently proved
upper bound gives an economical exact rank certificate. A modular
nonzero minor may also prove a characteristic-zero lower bound once
the reduction map and cleared denominators are justified; a modular
zero residual is insufficient for characteristic-zero containment.

The current-algebra part has more reusable support than a new
representation-theory library would suggest. In
\artifact{Mathlib/Algebra/Lie/Sl2.lean},
\leanid{IsSl2Triple.HasPrimitiveVectorWith} packages a highest-weight
vector. Its lemmas \leanid{lie_h_pow_toEnd_f},
\leanid{lie_e_pow_succ_toEnd_f}, and
\leanid{pow_toEnd_f_eq_zero_of_eq_nat} give the standard weight,
raising-coefficient, and finite-dimensional nilpotence calculations.
These can support the singlet polynomials and the bounded-weight
argument supplying \leanid{nilpotent_raising_saturation}.
They do not identify the present current actions with integrable
affine-VOA modules. Nor does the generic saturation theorem verify
the saved spanning matrix.

Mathlib already has formal vertex \emph{operators}:
\leanid{VertexOperator}, \leanid{VertexOperator.ncoeff}, and
\leanid{VertexOperator.of_coeff} are in
\artifact{Mathlib/Algebra/Vertex/VertexOperator.lean}.
They use Hahn-series linear maps from
\artifact{Mathlib/Algebra/Vertex/HVertexOperator.lean}, which provides
coefficient and composition operations. This is relevant groundwork
for mode indexing and lower truncation. The inspected file does not
provide a complete conformal vertex algebra with vacuum, translation,
Jacobi/locality, modules, and invariant forms; its own TODO list
includes locality and residue products. Directory and declaration
searches, followed by inspection of these files, located no existing
implementation of the manuscript's complete VOA interface. That is
a statement about the audited checkout, not a claim that Mathlib
contains no relevant vertex-algebra ingredients.

The filtered Wick criterion is therefore a worthwhile intermediate
deliverable. Its hypotheses should explicitly state the mode identities
and filtration controls used in the proof. Formalizing only an abstract
implication is useful, provided the final theorem continues to display
those hypotheses until the actual lattice construction supplies them.

For Appendix~\msref{sec:regularity}, the positive stress identity separates similarly.
Finite-grade self-adjoint operators and their nonnegative quadratic
forms can be treated with the inner-product-space library. For example,
\leanid{LinearMap.IsSymmetric.eigenvectorBasis} in
\artifact{Mathlib/Analysis/InnerProductSpace/Spectrum.lean} supports
finite-dimensional diagonalization. Turning the inequality into
regularity requires the specified conformal embeddings and the
module-theoretic finiteness results. A formal proof of a positive
matrix identity alone cannot establish those embeddings or
$C_2$-cofiniteness.

Appendix~\msref{sec:branching} has three distinct layers. The first is the core-top
enumeration and exact trace inequality; only the last linear implication
is currently checked. The second is finite algebra: branching
multiplicities, a based ring on four or six symbols, dimension equations,
and root-of-unity sums. These are credible conditional targets using
finite types, matrices, polynomial identities, and number theory.
The third is the interpretation through semisimple tensor categories,
formal codegrees, indicators, induction, and a modular center. It needs
structures and theorems beyond a table of fusion coefficients. In
particular, verifying a fusion table does not verify a pentagon,
construct its associator, or classify its categorifications.

The optional algebraic classification
Proposition~\msref{prop:optional-algebraic-identification}
(p.~\mspageref{prop:optional-algebraic-identification}) must have an explicit
parameter expressing exhaustion by the four categories in the cited
list. Its pivotal-scalar exclusion can then be formalized separately.
The repository's Evans--Gannon coefficient construction does not
discharge that exhaustion hypothesis. The manuscript's unitary route
avoids this optional assumption, but a Lean version of that route still
needs the separate unitary classification and Morita-equivalence
theorems. Avoiding one external premise does not automatically formalize
the replacement theorem.

For Appendix~\msref{sec:characters}, the installed library has \leanid{ModularForm} in
\artifact{Mathlib/NumberTheory/ModularForms/Basic.lean} and
\leanid{UpperHalfPlane.qExpansion},
\leanid{UpperHalfPlane.qExpansion_eq_zero_iff}, and
\leanid{UpperHalfPlane.qExpansion_coeff_unique} in
\artifact{Mathlib/NumberTheory/ModularForms/QExpansion.lean}.
There are also scalar theta, eta, congruence-subgroup, and level-one
developments. They do not directly supply Zhu's vector-valued
characters or the exact finite-image representation used here.
An appropriate first theorem would prove principal-part uniqueness
under explicit transformation, holomorphy, cusp, and invariant-vector
hypotheses, and then distinguish its application to the actual trace
functions from finite coefficient arithmetic.

\subsection{The analytic comparison in detail: A2, A3, and A1}

The target is Appendix~\msref{sec:analytic-unitarity},
\emph{\msnameref{sec:analytic-unitarity}}, beginning on
p.~\mspageref{sec:analytic-unitarity}.

The manuscript now gives mathematical proofs of all three comparison
steps. Their lack of a Lean formalization is not an assertion that they
remain mathematical conjectures. Conversely, replacing each of them
by a Lean proposition whose proof is an assumed field would only record
the dependency graph.

\paragraph{A2: the actual subnet and sector count.}
Proposition~\msref{clm:analytic-net-count}
(p.~\mspageref{clm:analytic-net-count}) is the precise target.
A local formalization can separate the branching multiplicities
$1,1,2$, the endomorphism algebra
$\mathbb C\oplus\mathbb C\oplus M_2(\mathbb C)$, and the rank-six
based-ring argument. Those are finite algebra once the representation
objects and their decompositions are supplied. The passage from an
algebraic VOA summand to a reducing Hilbert subspace requires the A3
common-core lemma, energy bounds, and an actual represented subnet.
The proof of local irreducibility uses locality, strong additivity,
the shared rotation action, and the separating vacuum; all four inputs
must be visible in a formal statement.

End-to-end A2 then needs finite-index local extensions, DHR sectors,
statistical dimension, the canonical algebra, center induction, and
complete rationality. The index-two exclusion also needs the precise
net-automorphism/VOA-automorphism correspondence and the finite-order
Lie-algebra fixed-rank argument. Its untwisted $\mathbb Z/3$ subcategory
comes from the transverse diagonal and antidiagonal isotropic lines,
followed by free-module induction. A formal proof should preserve that
construction: an abstract group of three invertible symbols does not
determine a trivial associator.

Mathlib does provide \leanid{VonNeumannAlgebra} and
\leanid{VonNeumannAlgebra.commutant} in
\artifact{Mathlib/Analysis/VonNeumannAlgebra/Basic.lean}.
Here a concrete von Neumann algebra is packaged by its bicommutant
property. This is useful support for locality and commutant arguments.
The file is not a conformal-net, DHR, subfactor-index, or Connes-fusion
library. Those theories were not located in the audited checkout.
The distinction matters: the basic operator-algebra vocabulary exists,
while the net-theoretic theorems used by A2 still need formalization.

\paragraph{A3: reconstruction, common cores, and exhaustion.}
The reconstruction statement is
Proposition~\msref{clm:analytic-reconstruction}
(p.~\mspageref{clm:analytic-reconstruction}), with its strong-unitarity
consequence in Proposition~\msref{prop:analytic-strong-unitarity}.
The most accessible analytic target is the graph-reduction
Lemma~\msref{lem:analytic-graph-reduction}
(p.~\mspageref{lem:analytic-graph-reduction}).
Represent a densely defined operator by \leanid{LinearPMap}, require
closability, and assume explicitly that the bounded orthogonal
projection preserves its domain and commutes there. The map
$(\xi,\eta)\mapsto(P\xi,P\eta)$ preserves the graph and therefore its
closure. This proves commutation on the closed domain. The adjoint
claim follows from the defining inner-product identity, using
$P^*=P$ and domain invariance. Both arguments can be formalized without
claiming that the original core is a core for the adjoint.

The existing declarations
\leanid{LinearPMap.IsClosable.graph_closure_eq_closure_graph},
\leanid{LinearPMap.IsClosable.closure_isClosed}, and
\leanid{LinearPMap.HasCore} are in
\artifact{Mathlib/Topology/Algebra/Module/LinearPMap.lean}.
The complementary file
\artifact{Mathlib/Analysis/InnerProductSpace/LinearPMap.lean} provides
\leanid{LinearPMap.IsFormalAdjoint} and
\leanid{LinearPMap.adjoint}.
Projection APIs include \leanid{Submodule.orthogonalProjectionOnto}
and \leanid{Submodule.starProjection} in
\artifact{Mathlib/Analysis/InnerProductSpace/Projection/Basic.lean}.
These are actual reusable declarations, not proposed abstractions.
The further conclusion about the von Neumann algebra generated by an
unbounded field requires an affiliation or polar/spectral-generation
interface in addition to the graph lemma.

The next target is the energy-domain argument. Starting with a positive
energy representation whose bounded-energy vectors form a weak module,
regularity gives an ordinary simple summand $S$. One must prove that
its Hilbert closure reduces $L_0$ and that its projection sends every
bounded-energy vector into a finite sum of finite-dimensional grades
of $S$. This establishes $PM\subset S\subset M$ before any field
commutation is used. Assuming finite multiplicities for the original
net representation would remove precisely the issue the proof resolves.
Likewise, a smeared field need not map the bounded-energy core to
itself; its value is a Hilbert-space vector. A Lean signature that
accidentally makes every smeared field an endomorphism of that core
would encode a stronger and inappropriate premise.

The grade-by-grade comparison of two positive invariant forms is another
bounded target. It first constructs an algebraic module endomorphism,
uses an eigenvector on a finite-dimensional lowest grade and simplicity
to prove that endomorphism scalar, and only then obtains boundedness of
the module isomorphism. This order avoids importing a Hilbert-space
Schur lemma before boundedness is known. Finite direct sums and the
counting argument that twelve reconstructed simple modules exhaust
twelve algebraic simples are comparatively small after the interfaces
are fixed.

Full A3 still requires the Henriques--Tener reconstruction theorem,
including its bounded-energy extension, equality of closed smeared
fields, and reconstruction of the original vacuum VOA. Formalizing
the consequences of an explicitly supplied reconstruction interface
would be useful. It would not be a formal proof of that theorem.

\paragraph{A1: pairwise annulus compatibility and positivity.}
The compatibility assertion is
Proposition~\msref{clm:analytic-annulus}
(p.~\mspageref{clm:analytic-annulus}); its domain and continuation steps
are Lemmas~\msref{lem:analytic-inverse-domain} and
\msref{lem:analytic-insertion-continuation}. The final application is
Theorem~\msref{thm:analytic-complete-unitarity}.
The revised proof deliberately constructs a lift for each compatible
annulus pair. It does not require a global canonical choice of lifts
for every annulus. A suitable formal interface would record the
following data and laws:
\begin{enumerate}
\item the centrally extended annulus semigroup and its localized
subsemigroups, together with their positive-energy representations;
\item the common localizer $\delta$, localized lifts $D_A,D_B$, and
the pair $\widetilde A=\widetilde\delta^{-1}D_A$,
$\widetilde B=D_B\widetilde\delta$, so that
$\widetilde B\widetilde A=D_BD_A$ is canonically localized;
\item the equality of local annulus actions with the represented
vacuum operators, incoming insertions, sewing, and the outgoing
unitary adjoint formula;
\item the round-annulus factorization giving both finite-energy
range inclusion and density of the transported form domain;
\item holomorphy on the sewn interior, coverage by all allowed
geometric representatives, and the no-insertion generators.
\end{enumerate}
The deck ambiguity must be retained: changing the common lift conjugates
the inserted field by $e^{2\pi i mL_0}$. Integer VOA weights make this
conjugation trivial even when module lowest weights are fractional.
Normalizing a vacuum operator alone does not eliminate the deck group
on those modules. This small algebraic cancellation is a realistic
formal target once the grading and representation laws are explicit.

The analytic continuation step can use
\leanid{AnalyticOnNhd.eqOn_of_preconnected_of_eventuallyEq} from
\artifact{Mathlib/Analysis/Analytic/Uniqueness.lean}, applied to scalar
matrix coefficients after all holomorphy and component hypotheses are
supplied. The finite-energy range and dense-domain lemmas can also be
isolated from annulus geometry. The final equality of concrete nets is
an elementary conditional commutant argument:
\[
 \mathcal A^{\mathrm{bl}}(I)\subset\mathcal N(I)
 \subset\mathcal A^{\mathrm{bl}}(I')'
 =\mathcal A^{\mathrm{bl}}(I).
\]
It explicitly uses locality and Haag duality on the same vacuum Hilbert
space; an abstract isomorphism of the two nets would not suffice.

These smaller arguments do not build the central extension, integrate
Virasoro representations, prove geometric sewing, or prove positivity
of Connes fusion and the resulting canonical VOA fusion pairings.
Those are the large missing foundations. An honest intermediate
formalization would name the exact integration, reconstruction, and
sewing hypotheses, prove their pairwise consequence, and keep the
positive-fusion theorem separate until it too is formalized.

\subsection{Categories, extensions, and the independent Evans--Gannon proof}

Mathlib has substantial general categorical infrastructure. Verified
examples include \leanid{CategoryTheory.BraidedCategory} in
\artifact{Mathlib/CategoryTheory/Monoidal/Braided/Basic.lean},
\leanid{CategoryTheory.RigidCategory} and duality interfaces in
\artifact{Mathlib/CategoryTheory/Monoidal/Rigid/Basic.lean}, and
\leanid{CategoryTheory.HalfBraiding}, \leanid{CategoryTheory.Center},
\leanid{CategoryTheory.Center.forget}, and
\leanid{CategoryTheory.Center.ofBraided} in
\artifact{Mathlib/CategoryTheory/Monoidal/Center.lean}.
These permit genuine reuse in Appendices~\msref{sec:branching},
\msref{sec:analytic-category}, and \msref{sec:analytic-full-theory}. They do not by
themselves provide a unitary fusion category with finite simple-object
data, positive spherical dimensions, modular $S,T$, a Q-system
correspondence, or the Haagerup classification. No such complete
manuscript-level framework was located in the local source audit.

A finite based ring can be formalized before its categorical realization.
One can then prove the conditional dimension equations and obstruction
calculations, keeping separate the theorem that a given unitary
category realizes that ring. Similarly, transporting an already
constructed algebra along a braided equivalence is distinct from
proving that it yields a local net or a unitary VOA extension.
The formal target for Appendix~\msref{sec:analytic-full-theory}
must preserve the manuscript's scope:
the original embedded extension and the chosen $\mathcal H_3$
extension need not be the same algebra; the 36 bimodule defects
preserve $U$; and the holomorphic-net construction is not identified
with the $E_8$ net by the VOA classification theorem.

The existing Evans--Gannon development should be credited precisely.
In \artifact{lean/HaagerupIzumi/Analytic/GammaConstruction.lean},
\leanid{HaagerupIzumi.GammaContinuation.hypGamma} constructs the
hyperbolic-gamma interface from the proved integral development.
\leanid{HaagerupIzumi.ruijsenaarsHypGamma} specializes it to the
required periods, and \leanid{HaagerupIzumi.EGEquations_ruijsenaars}
proves the equations for the resulting concrete coefficient data.
The theorem \leanid{HaagerupIzumi.exists_EGEquations_odd} assumes only
\leanid{Odd n} and produces complex coefficient data satisfying all
the packaged Evans--Gannon equations. The source genuinely supplies
the earlier analytic interface; this final theorem does not assume
an unconstructed gamma function.

The packaged proposition \leanid{HaagerupIzumi.CoeffData.EGEquations}
is defined in \artifact{lean/HaagerupIzumi/Completion.lean}; it records
the linear symmetries and axis sums together with the quadratic,
cubic, and quartic equations. The project README explicitly stops
before the reconstruction theorem and categorical conclusions of
the separate infinitude manuscript. Thus this development does not
prove exhaustion of rank-six fusion categories, construct the
present lattice subVOA, or establish its complete unitarity. Its
finite Fourier infrastructure can be reused: the local interpolation
proof already uses \leanid{ZMod.dft} and related declarations from
\artifact{Mathlib/Analysis/Fourier/ZMod.lean}. Its contour machinery
may be useful for appropriate scalar analytic sublemmas, but it is
not an annulus or conformal-net construction.

\subsection{A staged program with explicit completion criteria}

The following stages give reviewable deliverables without suggesting
that a short final Lean file would amount to a short foundational proof.
Some stages can proceed independently; their applications must retain
unproved inputs visibly.
\begin{enumerate}
\item \textbf{Certify finite data.} Define the exact coefficient algebra,
sparse vectors, and matrix witnesses; prove a small checker's soundness;
import and verify selected residual and rank certificates. The
deliverable is a theorem about named finite data, with the interpretation
map stated. It is not yet a theorem about the actual VOA.
\item \textbf{Complete the current-algebra layer.} Formalize the bounded
weight argument for joint nilpotence, the singlet projection, the
charged-module spanning deductions, and the top-table arithmetic.
Connect these to the already checked saturation and inequality lemmas.
The deliverable removes the corresponding finite algebraic assumptions.
\item \textbf{Formalize the recursive closure criterion.} Give a precise
graded mode-algebra interface and prove the filtered Wick and Jacobi
reductions. Instantiate the finite gates through the certified data.
Construction of the genuine lattice VOA and validation of that
instantiation remain explicit subsequent obligations.
\item \textbf{Separate representation theory from finite consequences.}
Formalize the core-top and branching deductions, the based-ring
classification, and character uniqueness under stated VOA and
categorical interfaces. Keep optional categorical exhaustion as a
separate parameter. The deliverable should make clear which deep
representation-theory theorems remain inputs.
\item \textbf{Formalize the analytic lemmas.} Prove common-core reduction,
adjoint commutation, dense-domain transport, grade-wise form comparison,
deck cancellation, and the conditional Haag-duality squeeze. State
reconstruction and annulus comparison with their full domains and
normalizations. These supply reusable analysis without pretending to
formalize the entire net/VOA correspondence.
\item \textbf{Discharge the foundational interfaces.} Build the required
VOA, module tensor, conformal-net, annulus, and unitary fusion theories;
formalize the cited reconstruction, classification, positivity, and
extension theorems; construct the actual objects and connect the finite
certificates. Only completion of these dependencies supports an
end-to-end Lean theorem matching the manuscript's full conclusion.
\end{enumerate}

\subsection{Assumptions, computation, and the trust boundary}

A theorem with explicit parameters such as a spanning identity or a
reconstruction theorem can be fully kernel-checked as an implication.
Its parameters are not new global logical axioms. Nevertheless, the
application to the CFT is unproved until those parameters are supplied.
An interface whose fields simply repeat the desired conclusion adds
little verification; useful interfaces expose the smaller mathematical
laws from which the conclusion is actually derived.

An admitted \leanid{sorry} introduces \leanid{sorryAx}; a declared
\leanid{axiom} installs an unproved proposition. Neither is equivalent
to a theorem with all premises discharged. A source search of the local
\leanid{HaagerupIzumi} Lean files found no uses of \leanid{sorry},
new \leanid{axiom} declarations, or \leanid{native_decide}; the saved
axiom reports give the more relevant dependency information for the
specific checked conclusions. A text search alone is not a substitute
for an axiom audit of the final theorem.

There is a version-specific point about computation. In this Lean
4.33.1 installation, \leanid{native_decide} calls
\leanid{Lean.Meta.nativeEqTrue}; the implementation in
\artifact{Lean/Meta/Native.lean} compiles and evaluates a closed Boolean
and, on success, introduces an auxiliary axiom asserting that its value
is true. The kernel then checks a proof relative to that axiom; it does
not independently replay the native computation by ordinary reduction.
The older \leanid{Lean.ofReduceBool} mechanism is deprecated in this
toolchain. Consequently a future certificate replay using
\leanid{native_decide} would need an explicit expanded trust statement,
even if it is useful operationally.

By contrast, ordinary proof-producing tactics and a \leanid{decide}
proof checked by kernel reduction can verify an exact certificate
without that native-evaluation axiom, subject to the final dependency
report. For large data, the practical target is a small sound checker
and manageable proof certificates, not necessarily replaying the
original search algorithm inside the kernel. A Boolean printed by
\texttt{\#eval}, a Python exact calculation, a successful hash check,
and a theorem with a kernel-checked proof are different artifacts.
They should remain separately labeled in the ledger.

Future compiler and certificate runs must retain the repository's
shared-desktop resource policy and run serially through its guarded
launcher. The preserved baseline assessment introduced no new Lean declarations and
does not claim any additional kernel-checked results beyond the
recorded files above.

\section{Additional assessment of the physics revision}
\label{sec:physics-lean-assessment}

The new physics sections provide much more explicit mathematical data:
character coefficients, Virasoro-primary multiplicities, spin-resolved
spectra, thermal intervals, low-weight couplings, rational correlators,
fusion coefficients, defects, and boundary partition functions. None
of the Python producers in \artifact{c8_physics_20260927/} is a Lean
proof or has a Lean-verified implementation. Their exact arithmetic
and internal checks improve reviewability and supply concrete inputs
for formalization; they do not enlarge the scopes of the five existing
CFT declarations listed in the preceding assessment. A separate new
file adds the two small thermal-comparison lemmas stated below, bringing
the selected CFT coverage to seven principal declarations in three
files. The five baseline declarations retain their original scope.

\subsection{Current manuscript anchors and historical concordance}

The linked numbers, titles, and pages in the following table refer to
the current manuscript. The first column is only a concordance for
readers of the archived 55-page PDF; its former section numbers are
not operative references in this guide. The current document places
the mathematical development in appendices after the physics sections.

\begin{longtable}{@{}P{0.11\linewidth}P{0.57\linewidth}P{0.24\linewidth}@{}}
\toprule
Archived label & Current manuscript title & Current anchor\\
\midrule\endfirsthead
\toprule Archived label & Current manuscript title & Current anchor\\
\midrule\endhead\bottomrule\endfoot
1 & \msnameref{sec:statement}
& Section~\msref{sec:statement}, p.~\mspageref{sec:statement}\\
2 & \msnameref{sec:construction}
& Appendix~\msref{sec:construction}, p.~\mspageref{sec:construction}\\
3 & \msnameref{sec:finite}
& Appendix~\msref{sec:finite}, p.~\mspageref{sec:finite}\\
4 & \msnameref{sec:gate4}
& Appendix~\msref{sec:gate4}, p.~\mspageref{sec:gate4}\\
5 & \msnameref{sec:gate5}
& Appendix~\msref{sec:gate5}, p.~\mspageref{sec:gate5}\\
6 & \msnameref{sec:properness}
& Appendix~\msref{sec:properness}, p.~\mspageref{sec:properness}\\
7 & \msnameref{sec:regularity}
& Appendix~\msref{sec:regularity}, p.~\mspageref{sec:regularity}\\
8 & \msnameref{sec:branching}
& Appendix~\msref{sec:branching}, p.~\mspageref{sec:branching}\\
9 & \msnameref{sec:characters}
& Appendix~\msref{sec:characters}, p.~\mspageref{sec:characters}\\
10 & \msnameref{sec:analytic-unitarity}
& Appendix~\msref{sec:analytic-unitarity}, p.~\mspageref{sec:analytic-unitarity}\\
11 & \msnameref{sec:analytic-category}
& Appendix~\msref{sec:analytic-category}, p.~\mspageref{sec:analytic-category}\\
12 & \msnameref{sec:analytic-full-theory}
& Appendix~\msref{sec:analytic-full-theory}, p.~\mspageref{sec:analytic-full-theory}\\
13 & \msnameref{sec:comparison}
& Appendix~\msref{sec:comparison}, p.~\mspageref{sec:comparison}\\
A & \msnameref{sec:ledger}
& Appendix~\msref{sec:ledger}, p.~\mspageref{sec:ledger}\\
B & \msnameref{sec:formalization}
& Appendix~\msref{sec:formalization}, p.~\mspageref{sec:formalization}\\
C & \msnameref{sec:compilation}
& Appendix~\msref{sec:compilation}, p.~\mspageref{sec:compilation}\\
\end{longtable}

The added physics material has the following current anchors.
``Focused'', ``substantial'', and ``foundational'' retain the
dependency-based meanings stated earlier.

\begin{longtable}{@{}P{0.70\linewidth}P{0.21\linewidth}@{}}
\toprule
Current section title & Current anchor\\
\midrule\endfirsthead
\toprule Current section title & Current anchor\\
\midrule\endhead\bottomrule\endfoot
\msnameref{sec:physics-spectrum}
& Section~\msref{sec:physics-spectrum}, p.~\mspageref{sec:physics-spectrum}\\
\msnameref{sec:physics-correlators}
& Section~\msref{sec:physics-correlators}, p.~\mspageref{sec:physics-correlators}\\
\msnameref{sec:physics-couplings}
& Section~\msref{sec:physics-couplings}, p.~\mspageref{sec:physics-couplings}\\
\msnameref{sec:physics-defects}
& Section~\msref{sec:physics-defects}, p.~\mspageref{sec:physics-defects}\\
\msnameref{sec:physics-H3}
& Section~\msref{sec:physics-H3}, p.~\mspageref{sec:physics-H3}\\
\msnameref{sec:physics-data-guide}
& Section~\msref{sec:physics-data-guide}, p.~\mspageref{sec:physics-data-guide}\\
\end{longtable}

\subsection{Characters, Virasoro primaries, and the full spectrum}

Section~\msref{sec:physics-spectrum},
\emph{\msnameref{sec:physics-spectrum}}, is in
\artifact{c8_physics_20260927/spectrum.tex}. Its producer
\artifact{spectrum.py} checks the published initial coefficients,
the rational matrix recurrence, nonnegative integrality through the
requested grade, descendant reconstruction, parity of the resulting
spin spectrum, and an independent $E_4/\eta^8$ identity. The default
data include 25 coefficients for each of twelve characters, complete
full-CFT spectra through $\Delta=24$, and an independent $E_8$ spectrum.
These are finite arithmetic conclusions conditional on the input
matrices and their character interpretation.

\begin{longtable}{@{}P{0.26\linewidth}P{0.64\linewidth}@{}}
\toprule
Physics argument & Missing Lean theorem and difficulty\\
\midrule\endfirsthead
\toprule Physics argument & Missing Lean theorem and difficulty\\
\midrule\endhead\bottomrule\endfoot
Matrix recurrence
& Define the twelve-dimensional exponent/characteristic data and a
finite rational recurrence, prove nonzero denominators, and prove that
its finite output agrees with the formal differential equation.
This is focused to substantial. Identifying the resulting vector with
the actual VOA trace functions still requires the foundational
character theorem assessed under Appendix~\msref{sec:characters}.\\
Virasoro-primary subtraction
& The finite Euler-product convolution and its inverse are focused.
Their interpretation needs the irreducible $c=8$ Virasoro characters,
the Kac determinant or an equivalent irreducibility theorem, and the
unitary highest-weight decomposition of each module. These are
substantial to foundational representation-theoretic inputs.\\
Complete spin-resolved cutoff
& Given nonnegative lowest weights and coefficient functions, prove
the product formula and that no grade above $N$ contributes to
$\Delta\leq N$. The finite sums, parity identity, and relevant/marginal
counts are focused. The tensor-product Hilbert-space interpretation
and the diagonal modular invariant are additional premises.\\
$E_8$ comparison and invariant counts
& A finite identity between the oscillator series and the selected
character combination is focused. It does not identify the extension
algebra. Counts $1488=6\cdot248$ for one chiral Haagerup action and
$36=6^2$ for both independent actions additionally require the precise
invariant-space theorem. That categorical interpretation belongs to
the separate defect formalization below.\\
\end{longtable}

There is useful existing power-series infrastructure.
\leanid{PowerSeries.coeff_mk} and \leanid{PowerSeries.coeff_mul} occur
in \artifact{Mathlib/RingTheory/PowerSeries/Basic.lean}.
The file
\artifact{Mathlib/Combinatorics/Enumerative/Pentagonal/PowerSeries.lean}
contains the Euler-product development, including
\leanid{PowerSeries.coeff_prod_one_sub_X_pow_eventually_eq} and
\leanid{PowerSeries.WithPiTopology.hasProd_one_sub_X_pow}.
These declarations were read in the pinned checkout. They support
formal coefficient manipulations; they do not state that these
coefficients count this CFT's Virasoro primaries.

An economical formalization would first take an abstract sequence of
graded dimensions and prove the descendant-removal and tensor-product
counting lemmas. A second layer would connect that sequence to the
Evans--Gannon recurrence. The final interpretation theorem must connect
both layers to the actual modules. Merely importing the printed
integers as definitions and proving that their sum is $2941$ would
check arithmetic but leave the substantial spectral assertion open.

\subsection{Thermal bounds and their numerical intervals}

Proposition~\msref{prop:physics-modular-tail}
(p.~\mspageref{prop:physics-modular-tail}) proves the modular heavy-tail
estimate in Section~\msref{sec:physics-spectrum}. For
$E=\Delta-2/3$, an exact cutoff $D>2/3$, and $b'=4\pi^2/b$, it uses
\[
 Z=L_D+H_D,\qquad Z(b)=Z(b'),\qquad
 H_D(b)\leq e^{-(b-b')(D-2/3)}H_D(b')
\]
to bound the tail at $b>2\pi$. Two real-number steps are now checked in
\artifact{lean/HaagerupIzumi/CFT/PhysicsBounds.lean}, in namespace
\leanid{HaagerupIzumi.CFT}:
\begin{itemize}
\item \leanid{modular_heavy_tail_comparison} assumes
$L+H=L_d+H_d$, $H\leq rH_d$, and $0\leq r<1$, and proves
$H\leq r(L-L_d)/(1-r)$;
\item \leanid{selfdual_tail_first_moment_comparison} assumes
$E_{\rm cut}>0$ and $E_{\rm cut}H\leq L'$, and proves
$H\leq L'/E_{\rm cut}$.
\end{itemize}
The successful guarded compiler record is
\artifact{c8_physics_20260927/physics_lean.log}, with sampled peak
group RSS $1.549$~GiB. Both axiom reports contain only
\leanid{propext}, \leanid{Classical.choice}, and \leanid{Quot.sound};
the source contains neither \texttt{sorry} nor new axioms. The
nonnegative contraction range is retained explicitly, although the
first rearrangement only needs $r<1$ once its comparison is supplied.

These lemmas do not derive the comparison from a spectral sum or
identify $r$ with the exponential in the application. The actual
infinite-sum theorem remains more substantial: it needs convergence,
termwise positivity, completeness of the finite cutoff, and the
modular functional equation. Its self-dual form also needs
differentiability and the identity $Z'(2\pi)=0$. The first and second
moment bounds require justified moment convergence and the elementary
optimization of $E^ke^{-aE}$ on a positive half-line.

The Python output is a conditional interval certificate. Its positive
spectrum and modular invariance are mathematical premises, not facts
established by evaluating finitely many coefficients. The implementation
uses directed decimal arithmetic, a fixed rational enclosure of $\pi$,
and correctly rounded exponential/logarithm functions; the displayed
six-decimal table is accepted only when every rigorous interval lies
within $10^{-6}$ of its printed value. None of that implementation is
currently kernel checked.

For a Lean version, a small rational interval checker with proved
exponential and logarithm enclosures would be preferable to treating
Python's decimal library as an axiom. One can separately formalize the
analytic tail theorem over the real numbers and feed it certified
finite moment bounds. Existing real analysis and infinite-sum
infrastructure is relevant, but the complete pipeline from characters
to the saved numerical intervals is new work. The two checked
implications are the completed local algebraic part; sound interval
evaluation and its analytic interpretation
are substantial. The CFT assumptions still inherit the foundational
requirements of the earlier sections.

\subsection{Currents, lattice correlators, and low-weight couplings}

The correlator and coupling material in
Sections~\msref{sec:physics-correlators} and
\msref{sec:physics-couplings} is supplied by
\artifact{c8_physics_20260927/correlators.tex} and
\artifact{c8_physics_20260927/couplings.tex}. Its data distinguish
direct finite lattice-mode products, rational identities reconstructed
from saved singular OPEs, and coefficients completed using current
Ward identities. Those three forms of evidence have different
interpretation obligations; the final reports specify the supplied
bases and direct probes.

In particular the completed weight-two dataset in
\artifact{c8_physics_20260927/full_weight_two/} uses all $113$ chiral
Virasoro primaries. Its saved report records $1{,}442{,}897$ ordered
tensor slots, $18{,}365$ nonzero entries, $51{,}076$ current Ward
checks, and $33$ additional direct lattice-mode probes beyond the
seed calculation. This supplies substantially more explicit input
than a list of summed OPE coefficients. Those successful Python
checks remain separate from the Lean obligations below.

The following finite targets are realistic once the coefficient
algebra and basis conventions are represented in Lean:
\begin{itemize}
\item singular current brackets, levels, gauge factors, pairings, and
the orthogonality and positivity of the stated finite Gram matrices;
\item polynomial identities obtained by clearing denominators in
the displayed rational four-point functions, with collision loci
excluded when interpreting them as functions;
\item exact extraction of three-point coefficients and their squared
unit-norm values, retaining the specified invariant/Hermitian pairing
and adjoint conventions;
\item the finite Ward recursion, its compatibility with both current
factors, and uniqueness of its completion from the supplied seed slice;
\item the completeness of the chosen low-weight primary subspace,
provided its exact rank and descendant-complement certificates.
\end{itemize}
These are substantial finite algebra rather than a need to formalize
all CFT at once. The bridge from encoded lattice words to actual mode
operators is the same missing bridge described for the finite gates.
A successful check of a completed tensor against Ward identities
establishes consistency of the completion. Its identification with
all actual mode products also needs a proved determination theorem
and authenticated seed products; a finite selection of direct probes
does not by itself check every omitted product.

Beyond that finite layer, one needs the state-field correspondence,
locality, conformal covariance, OPE convergence or formal expansion
identities, and an explanation of how the chiral functions assemble
into local bulk correlators. The known Potts and parafermion inputs
also need their formal embeddings. These are foundational. A chiral
vacuum correlator or a chosen lattice intertwiner block is not by
itself the four-point function of the lightest diagonal bulk scalar.
The formal statement must preserve the external fields and their
normalization rather than silently promote one kind of block to the
other.

No claim that a dimension-two scalar is exactly marginal follows
from its existence or its three-point coefficient alone. A formal
statement about a conformal manifold would additionally require a
precise deformation theory, beta functions, and the appropriate
all-orders argument. That subject is outside both the present Lean
files and the finite arithmetic producer.

\subsection{Modular data, defects, and conformal boundaries}

Sections~\msref{sec:physics-defects} and \msref{sec:physics-H3},
\emph{\msnameref{sec:physics-defects}} and
\emph{\msnameref{sec:physics-H3}}, are in
\artifact{c8_physics_20260927/defects.tex}. Their
producer \artifact{defects.py} computes in
$\mathbb Q(\zeta_{13})$, with exact Verlinde coefficients, modular
relations, the forgetful multiplicities, and boundary/defect spectral
tables. These checks establish finite algebraic properties of the
input matrices and based rings, not the existence of the required
associator or a conformal defect.

\begin{longtable}{@{}P{0.27\linewidth}P{0.63\linewidth}@{}}
\toprule
Physics argument & Missing Lean theorem and difficulty\\
\midrule\endfirsthead
\toprule Physics argument & Missing Lean theorem and difficulty\\
\midrule\endhead\bottomrule\endfoot
Cyclotomic modular arithmetic
& Prove that the twelve-coordinate quotient representation is a
field embedded in $\mathbb C$, fix the primitive root and conjugation,
and verify the Gauss-sum square root and sign. Then $S^2=1$, the
twist relation, Verlinde integrality, and fusion associativity are
substantial but bounded exact computations.\\
Diagonal Verlinde lines
& Given a unitary modular tensor category and the diagonal full CFT,
derive line eigenvalues, twisted character matrices, and their
sector multiplicities. The scalar sums are focused; the theorem that
these matrices describe actual topological defects requires the
categorical/full-CFT construction and is foundational.\\
Cardy boundaries and annuli
& Formalize the boundary-state normalization, open/closed-channel
identity, and the positive square roots entering $g$ factors.
Finite character multiplicities are focused. Their interpretation
as boundary Hilbert spaces requires the Cardy/extension theorem.
Formal verification of annulus integers is not that theorem.\\
Selected chiral $\mathcal H_3$ action
& Prove the center-forgetful correspondence for the chosen extension,
the bimodule/full-defect equivalence, and the identification of tube
multiplicity spaces with $M_\mu\otimes\overline{H'}$.
These are foundational categorical inputs. Once supplied, the
$6\times12$ forgetful ring-map check and charge matrices are
substantial finite algebra.\\
Invariant and charged gaps
& Under the correct action, derive the fixed spaces and count their
lowest scalar grades. This is focused after the sector assignment.
It must distinguish one chiral action from both independent factors,
and tube multiplicity spaces from their multiple copies in actual
defect Hilbert spaces. Modular covariance alone is insufficient.\\
Six extension boundaries
& Verify the finite annulus coefficients from the forgetful functor
and the six-object fusion ring. The data-to-category interface must
retain the actual module category, duals, and boundary normalization.
The finite arithmetic is substantial; constructing these boundaries
as part of the full CFT remains foundational.\\
\end{longtable}

The pinned checkout defines \leanid{CyclotomicField} and
\leanid{IsCyclotomicExtension} in
\artifact{Mathlib/NumberTheory/Cyclotomic/Basic.lean}; it also has
primitive-root, complex-root-of-unity, and cyclotomic-polynomial
developments. Thus the field-theory layer need not start from scratch.
It still takes work to prove that the producer's concrete coordinate
arithmetic and chosen complex embedding agree with those structures.
An arbitrary root of the polynomial without the intended complex
embedding does not determine the positive quantum dimensions used
by the physics interpretation.

\subsection{Bootstrap placement and observable scope}

The physics reading guide in Section~\msref{sec:physics-data-guide}
distinguishes the diagonal full theory from
the selected holomorphic extension and states which published bootstrap
plots use which symmetry sectors. Formalizing the label correspondence,
the central-charge convention $c_L=c_R=8$, and a rational gap extracted
from a proved spectrum would be focused. It would not formally verify
a published numerical exclusion curve. That requires the original
functional, its positivity on every required conformal block, the
truncation/error bounds, and a sound numerical or exact certificate.
A plotted point read from a figure is evidence for comparison, not a
kernel-checked exclusion or inclusion theorem.

The observable inventory is also an interface specification. Characters
determine unflavoured multiplicities and, with analytic control,
partition functions. They do not determine every normalized OPE
coefficient or defect eigenvector. The new coupling and defect datasets
close specific parts of that gap; each needs its own sound checker and
interpretation theorem. Keeping those claims separate is essential to
an eventual formal treatment that a physicist can use without reading
the implementation.

\subsection{A concrete order of formalization}

A useful first deliverable is a small exact coefficient library covering
the rational recurrence, finite Euler product, primary reconstruction,
and cutoff-completeness lemma. Its theorem should expose the input
character interpretation explicitly. A second deliverable is a sound
finite algebra checker for the cyclotomic modular data and the supplied
low-weight tensors, with an explicit interpretation map. A third is
the real-analytic tail theorem and a certified interval evaluator.
Each can be independently reviewed before the foundational VOA and
defect interfaces are complete.

Completion criteria must be semantic: the formal theorem should refer
to the actual character coefficients, mode products, or represented
sectors through proved maps. File hashes, successful Python assertions,
or a Lean theorem about unrelated constants cannot replace that
connection. The two explicitly stated thermal-comparison lemmas are
the additional Lean results of this physics revision; the much larger
analytic and physical interpretations remain to be formalized.
